\begin{textsample}{POS Dim 3 – ba – Score 13.00 – ba/ba\_ef\_11.txt}  \label{ex:f3_pos_180}
8.2.3.
Educação Física O processo de construção da proposta curricular da Educação Física Escolar para a rede pública e privada da Bahia foi permeado pela compreensão do contexto da legalidade e legitimidade da BNCC, das Orientações Curriculares do Ensino Fundamental de nove anos25, da dinâmica estabelecida pela \textbf{instituição} responsável por esse processo ( Secretaria de Educação e seus atores sociais ), das \textbf{sugestões} elencadas pelo processo de consulta pública e da produção científica da área.
Sendo assim, pensar o currículo da Educação Física Escolar ( EFE ) constitui -se, pedagogicamente, como possibilidade de construção 25 BAHIA.
Secretaria da Educação.
Superintendência de Desenvolvimento da Educação Básica.
Diretoria de Educação Básica.
Orientações curriculares e subsídios didáticos para a organização do trabalho pedagógico no ensino fundamental de nove anos.
Salvador, 2013. da cidadania com autonomia intelectual, ética e moral, por meio dos conhecimentos historicamente construídos e fundamentados legalmente neste componente curricular obrigatório da Educação Básica, integrado à proposta pedagógica da escola ( BRASIL, 2003 ).
Adicionalmente, compreende -se que este componente deverá ser ofertado em todos os níveis e modalidades de ensino, ministrado por docentes licenciados em Educação Física.
Nesse cenário, esclarecemos que a efetividade do currículo dependerá das condições objetivas de implantação e execução em cada escola baiana, considerando elementos sociais, políticos, econômicos, pedagógicos, didáticos e formativos do processo, reconhecendo a individualidade de cada realidade escolar.
De outra forma, identifica -se um elemento primordial para a efetivação e vida do currículo na rede, o processo de formação continuada e permanente, para que as ações pedagógicas do DCRB sejam efetivadas e possam interferir na realidade social dos escolares.
Como tal, a Educação Física Escolar ( EFE ), no contexto da Área das Linguagens, configura -se como relevante para o processo de formação e desenvolvimento integral dos estudantes, durante o Ensino Fundamental, oferecendo possibilidades enriquecedoras de ampliação cultural do potencial dos escolares de intervirem de maneira crítica, autônoma e criativa na realidade social, por meio da pluralidade das práticas corporais sistematizadas e das representações sociais26.
Assim, a ampliação cultural aqui referendada compreende saberes e práticas corporais, experiências estéticas, emotivas e \textbf{lúdicas}, que se inscrevem, mas não se restringem à racionalidade típica dos saberes científicos que, comumente, orientam as práticas pedagógicas na escola.
Além disso, as experiências irrestritas com as práticas corporais e a segurança que esse conhecimento pode oferecer a cada estudante lhe oportunizarão experiências de autonomia e segurança em contextos de saúde e lazer, que, na vida do ser humano trabalhador moderno, tomam contornos ainda mais relevantes e fundamentais.
Para pensar uma proposta de currículo, é preciso reconhecer que os estudantes do Ensino Fundamental possuem modos próprios de vida e múltiplas experiências pessoais e sociais, o que torna necessário reconhecer a existência de \textbf{infâncias}, no plural, e, consequentemente, a singularidade de qualquer processo escolar e sua interdependência com as características da comunidade local. 26 GONZÁLEZ, F.
J. ; FRAGA, A.
B.
Afazeres da Educação Física na escola : planejar, ensinar, partilhar.
Erechim : Edelbra, 2012.
As \textbf{crianças} possuem conhecimentos que precisam ser, por um lado, reconhecidos e problematizados nas vivências escolares, com vistas a proporcionar a compreensão do mundo e, por outro, ampliados de maneira a potencializar a inserção e o trânsito dessas \textbf{crianças} nas várias esferas da vida social ( BRASIL, 2017 ).
Diante do compromisso com a formação estética, sensível e ética, a Educação Física, aliada aos demais componentes curriculares, assume, nessa composição curricular, o papel com a qualificação para a leitura, a produção e a vivência das práticas corporais.
Para tanto, entende -se que os professores devem buscar formas de trabalho pedagógico pautadas no diálogo, considerando a impossibilidade de ações uniformes ou lineares, que possam atender às demandas específicas de grupos naturalmente não incluídos.
Nas aulas, as práticas corporais poderão ser compreendidas como fenômeno cultural dinâmico, diversificado, pluridimensional, singular e contraditório.
Desse modo, é possível assegurar aos estudantes a (re)construção de um conjunto de conhecimentos que permitam ampliar sua consciência a respeito de seus movimentos e dos recursos para o \textbf{cuidado} de si e dos outros, além de desenvolver autonomia para apropriação e utilização da cultura corporal em diversas finalidades humanas, favorecendo sua participação de forma confiante e autoral na sociedade ( BRASIL, 2017 ).
Sendo assim, considera -se que haverá ampliação do acervo cultural corporal dos estudantes do Ensino Fundamental se todos os conhecimentos tematizados contemplarem a inclusão como princípio de suas ações pedagógicas, de modo que tanto o público ora excluído quanto aqueles que não demandam tratamento específico desfrutem das aprendizagens desejadas para essa etapa educacional.
No panorama da Base Nacional Comum Curricular ( BNCC ), a Educação Física Escolar é compreendida como fenômeno cultural dinâmico, diversificado, pluridimensional, singular e contraditório, tematizada por meio das práticas corporais em suas diversas formas de “ codificação e significação social, entendidas como manifestações das possibilidades expressivas dos sujeitos, produzidas por diversos grupos sociais no decorrer da história ” ( BRASIL, 2017, p. 217 ).
Neste documento, compreende -se que há três elementos fundamentais comuns às práticas corporais : “ movimento corporal como elemento essencial ; organização interna ( de maior ou \textbf{menor} grau ), pautada por uma lógica específica ; e produto cultural vinculado com o lazer/entretenimento e/ou o \textbf{cuidado} com o corpo e a saúde ” ( BRASIL, 2017, p. 211 ).
Vale ressaltar que a conceituação de práticas corporais necessita atender a esses três elementos fundamentais, além de serem aquelas realizadas fora das obrigações laborais, \textbf{domésticas}, higiênicas e religiosas, nas quais os sujeitos se envolvem em função de propósitos específicos, sem caráter instrumental.
Essa condição assertiva evita que qualquer movimento corporal seja inserido no currículo sem critério ou relação direta com o intento pedagógico do componente no Ensino Fundamental.
Dito isso, a Educação Física no Ensino Fundamental oferecerá, por meio das práticas corporais sistematizadas e das possibilidades de se movimentar, acesso a uma dimensão de conhecimentos e de experiências complementadas dentro e fora do ambiente escolar.
Considerando a extensão geográfica e a multiplicidade cultural que compõem a identidade do Estado da Bahia, a proposta curricular em tela leva em consideração a organização geográfica dos 27 Territórios de Identidade que marcam o Estado da Bahia, pois necessitam ser considerados nas proposições curriculares de cada território, cidade e escola.
Desta forma, diante da proposta indicada pela BNCC e das Orientações Curriculares do Ensino Fundamental de 9 ( nove ) anos na Bahia, o desenho curricular proposto para Educação Física Escolar organizará o conhecimento e as unidades temáticas sustentadas nas discussões de González e Schwengber27 e de González e Fraga28, a saber : a ) possibilidades do se-movimentar : abordadas como oportunidades de ampliação dos conhecimentos do próprio corpo, em diversos espaços e tempos em múltiplos contextos culturais.
Sendo assim, EFE oportunizará às \textbf{crianças} desafios psicomotores e cognitivos, na construção de novas referências sobre seu próprio corpo, de potencialidades para se-movimentar e de interação com o ambiente e com outros.
Além disso, destacamos que apesar da experiência de movimento ocupar um lugar central, não acontecerá no vazio social, pois estará permeada de valores e formas de entender o mundo. b ) a segunda dimensão do conhecimento da EFE se refere ao estudo das práticas corporais sistematizadas, com alguns elementos em comum, como : 1 ) o movimento corporal como elemento essencial ; 2 ) uma organização interna ( de maior ou \textbf{menor} grau ) pautada por uma lógica específica ; e 3 ) serem produtos culturais vinculados com o lazer/entretenimento e/ou o \textbf{cuidado} do corpo e a saúde.
Nessa perspectiva, as práticas corporais que fazem parte do campo de estudo da EF são : as acrobacias, as atividades aquáticas, as danças, os esportes, os exercícios físicos, os jogos e \textbf{brincadeiras}, as lutas, as práticas corporais de aventura na natureza, as ginásticas, a capoeira, a saúde e o lazer e práticas corporais. 27 GONZÁLEZ, F.
J. ; SCHWENGBER, M.
S.
V.
Práticas pedagógicas em Educação Física : espaço, tempo e corporeidade.
Erechim : Edelbra, 2012. 28 GONZÁLEZ, F.
J. ; FRAGA, A.
B.
Afazeres da Educação Física na escola : planejar, ensinar, partilhar.
Erechim : Edelbra, 2012. culturais vinculados com o lazer/entretenimento e/ou o \textbf{cuidado} do corpo e a saúde.
Nessa perspectiva, as práticas corporais que fazem parte do campo de estudo da EF são : as acrobacias, as atividades aquáticas, as danças, os esportes, os exercícios físicos, os jogos e \textbf{brincadeiras}, as lutas, as práticas corporais de aventura na natureza, as ginásticas, a capoeira, a saúde e o lazer e práticas corporais. c ) representações sociais sobre os conhecimentos da cultura corporal de movimento : entendidas como conhecimentos sociais construídos no campo científico, embasados na sociologia, antropologia, política, saúde coletiva, epidemiologia, fisiologia e anatomia, que contribuirão na formação humana.
Nesse contexto, a EFE problematizará conceitos sobre a origem e a dinâmica de transformação nas representações e práticas que se relacionam com as atividades corporais de tempo livre, o \textbf{cuidado} e a educação do corpo, seus vínculos com a organização da vida coletiva e individual, bem como os agentes sociais envolvidos em sua produção, tais como : o Estado, o mercado, a mídia, as \textbf{instituições} esportivas, as organizações sociais, as questões de gênero, socioeconômicas, políticas etc.
Ao \textbf{brincar}, dançar, jogar, praticar esportes, ginásticas ou atividades de aventura, para além da ludicidade, os estudantes se apropriam das lógicas intrínsecas ( regras, códigos, rituais, sistemáticas de funcionamento, organização, táticas etc. ) a essas manifestações, assim como trocam entre si e com a sociedade as representações e os significados que lhes são atribuídos.
Sendo assim, a Educação Física Escolar no Ensino Fundamental será um campo de experiências de se movimentar, das práticas corporais sistematizadas e das representações sociais, dispostas conforme o quadro a seguir “ Organizador Curricular ”, o qual discorre sobre a distribuição das unidades temáticas, competências específicas, objetos de conhecimento e habilidades, distribuídas em ciclos.
Por outro lado, a Figura 6, apresentada a seguir, representa as possibilidades que poderão ser adequadas a cada realidade escolar.
A proposta curricular deste documento sugere que os conhecimentos da EFE, delimitados em habilidades que privilegiem oito dimensões de conhecimento ( BRASIL, 2017 ), nas quais utilizaremos os exemplos que envolvem casos de inclusão e busquem facilitar o entendimento do docente : Experimentação : refere -se à dimensão do conhecimento que se origina pela vivência das práticas corporais, pelo envolvimento corporal na realização dessas práticas.
São conhecimentos que não podem ser acessados sem passar pela vivência corporal, sem que sejam efetivamente experimentados.
Trata -se de uma possibilidade única de apreender as manifestações culturais tematizadas pela Educação Física e de o estudante se perceber como sujeito “ de carne e osso ”.
Faz parte dessa dimensão, além do imprescindível acesso à experiência, \textbf{cuidar} para que as sensações geradas no momento da realização de uma determinada vivência sejam positivas ou, pelo menos, não sejam desagradáveis a ponto de gerar rejeição à prática em si.
O caminhar de um cego nunca será um conhecimento efetivo até que o estudante vidente seja desafiado a experimentar as suas tarefas naturalizadas do \textbf{dia} a \textbf{dia} de olhos vendados, com e sem companhia.
Não há outra forma de aprender esse conhecimento se não for através da vivência. - Uso e apropriação : refere -se ao conhecimento que possibilita ao estudante ter condições de realizar de forma autônoma uma determinada prática corporal.
Trata -se do mesmo tipo de conhecimento gerado pela experimentação ( saber-fazer ), mas dele se diferencia por possibilitar ao estudante a competência necessária para potencializar o seu envolvimento com práticas corporais no lazer ou para a saúde.
Diz respeito àquele rol de conhecimentos que viabilizam a prática efetiva das manifestações da cultura corporal de movimento não só durante as aulas, como também para além delas.
Ainda utilizando a cegueira como referência explicativa, a partir da experimentação, o estudante poderá e deverá desenvolver melhor a utilização de seus sentidos táteis e auditivos para usar e efetivamente se apropriar desse conhecimento. - Fruição : implica a apreciação estética das experiências sensíveis geradas pelas vivências corporais, bem como das diferentes práticas corporais oriundas das mais diversas épocas, lugares e grupos.
Essa dimensão está vinculada à apropriação de um conjunto de conhecimentos que permita ao estudante desfrutar e ser competente em uma prática corporal, de poder dar conta das exigências colocadas no momento de sua realização no contexto do lazer.
Trata -se de um grau de domínio da prática que permite ao sujeito uma atuação que lhe produz satisfação.
Ao incorporar em sua vida conhecimentos que podem lhe ser úteis no \textbf{dia} a \textbf{dia}, o estudante também apresentará a condição de fruição acerca desse conhecimento, na medida em que se enxergará competente e mais seguro no caso de a cegueira acometer alguém da família ou de lidar com esse público na vida em sociedade. - Reflexão sobre a ação : refere -se aos conhecimentos originados na observação e na análise das próprias vivências corporais e daquelas realizadas por outros.
Vai além da reflexão espontânea gerada em toda experiência corporal.
Trata -se de um ato intencional, orientado a formular e empregar estratégias de observação e análise para : ( a ) resolver desafios peculiares à prática realizada ; ( b ) apreender novas modalidades ; e ( c ) adequar as práticas aos interesses e às possibilidades próprios e aos das pessoas com quem compartilha a sua realização.
Sendo ou não um cego, a experiência com esses conhecimentos lhe permitirá refletir sobre as condições sociais que envolvem esse público e sua vida cotidiana, tornando -se, certamente, um agente efetivo na luta por condições melhores de vida social. - Construção de valores : vincula -se aos conhecimentos originados em discussões e vivências no contexto da tematização das práticas corporais que possibilitam a aprendizagem de valores e normas voltadas ao exercício da cidadania em prol de uma sociedade democrática.
A produção e \textbf{partilha} de atitudes, normas e valores ( positivos e negativos ) são inerentes a qualquer processo de socialização.
No entanto, essa dimensão está diretamente associada ao ato intencional de ensino e de aprendizagem e, portanto, demanda intervenção pedagógica orientada para tal fim.
Por esse motivo, deve -se focar a construção de valores relativos ao respeito às diferenças e ao combate aos preconceitos de qualquer natureza.
Ainda assim, não se pretende propor o tratamento apenas desses valores, ou fazê -lo só em determinadas etapas do componente, mas assegurar a superação de estereótipos e preconceitos expressos nas práticas corporais.
É muito mais simples não se preocupar com a inclusão quando se está incluído.
Essa reflexão determina um perfil de aprendizado em que valores verdadeiros da vida em sociedade se colocarão como desafio para a vida dos estudantes, haja vista que, ao experimentarem e refletirem sobre a cegueira, eles estarão reconstruindo seus valores e colocando em tela novos desafios para a construção de uma sociedade mais justa e em melhores condições de igualdade. - Análise : está associada aos conceitos necessários para entender as características e o funcionamento das práticas corporais ( saber sobre ).
Essa dimensão reúne conhecimentos como a classificação dos esportes, os sistemas táticos de uma modalidade, o efeito de determinado exercício físico no desenvolvimento de uma capacidade física, entre outros.
Essa dimensão do conhecimento permitirá ao estudante adentrar o mundo paralímpico, em suas regras, normas e modos de pensar o esporte e as práticas corporais, a ponto de compreender conceitos como “ classificação funcional ” e perceber sua relevância e interferência na prática esportiva profissional paralímpica. - Compreensão : está também associada ao conhecimento conceitual, mas, diferentemente da dimensão anterior, refere -se ao esclarecimento do processo de inserção das práticas corporais no contexto sociocultural, reunindo saberes que possibilitam compreender o lugar das práticas corporais no mundo.
Em linhas gerais, essa dimensão está relacionada a temas que permitem aos estudantes interpretar as manifestações da cultura corporal em relação às dimensões éticas e estéticas, à época e à sociedade que as gerou e as modificou, às razões da sua produção e transformação e à vinculação local, nacional e global.
Por exemplo, pelo estudo das condições que permitem o surgimento de uma determinada prática corporal, em uma dada região e época, ou os motivos pelos quais os esportes praticados por homens têm uma visibilidade e um tratamento midiático diferente dos esportes praticados por mulheres.
Ou estudos que mostrem os estereótipos construídos acerca da inutilidade de uma pessoa com deficiência podem oferecer elementos concretos que, imbricados nas dimensões anteriores, fortalecem a condição de esclarecimento acerca dos contextos socioculturais em que vivem. - Protagonismo comunitário : refere -se às atitudes/ações e conhecimentos necessários para os estudantes participarem de forma confiante e autoral de decisões e ações orientadas a democratizar o acesso das pessoas às práticas corporais, tomando como referência valores favoráveis à convivência social.
Contempla a reflexão sobre as possibilidades que eles e a comunidade têm ( ou não ) de acessar uma determinada prática no lugar em que moram, os recursos disponíveis ( públicos e privados ) para tal, os agentes envolvidos nessa configuração, entre outros, bem como as iniciativas que se dirigem para ambientes além da sala de aula, orientadas a interferir no contexto em busca da materialização dos direitos sociais vinculados a esse universo.
De posse desse perfil de conhecimentos, certamente pode -se criar a expectativa de que o estudante com essa formação se tornará um agente protagonista das ações em sua comunidade de moradia ou na comunidade de trabalho ou estudo.
Vale ressaltar que não há nenhuma hierarquia entre essas dimensões, tampouco uma ordem necessária para o desenvolvimento do trabalho no âmbito didático.
Cada uma delas exige diferentes abordagens e graus de complexidade para que se tornem relevantes e significativas.
Porém, é fundamental que cada uma dessas dimensões seja referência para o trabalho pedagógico, e, por consequência, seja o ponto de diálogo com os processos avaliativos dos docentes da Educação Física, de modo que estes tenham bem claro o que estão ensinando e, portanto, o que devem verificar na aprendizagem dos estudantes no processo avaliativo.
Na organização curricular, as unidades temáticas estão articuladas, pedagogicamente, considerando as características dos conhecimentos acumulados da Educação Física, dos professores, do contexto social e cultural da escola, dos alunos e alunas atreladas às competências gerais e específicas do componente curricular e das habilidades propostas do quadro organizador.
Além disso, a escola e o docente devem considerar esses pressupostos e observar a articulação com as competências gerais da BNCC e as competências específicas da Área de Linguagens, de modo que o componente curricular de Educação Física possa garantir aos estudantes o desenvolvimento de competências específicas ao final de seu ciclo de Ensino Fundamental.
A busca pelo desenvolvimento dessas 10 competências específicas da Educação Física, ao final do Ensino Fundamental, será definida pela orientação a partir de Unidades Temáticas nas quais estarão elencadas diversas habilidades a serem desenvolvidas pelos docentes.
As proposições temáticas poderão ser ampliadas a partir das experiências dos professores e das professoras, das características da realidade local, dos avanços da produção científica da área, das tecnologias disponíveis, bem como por meio da articulação com outras áreas do conhecimento, considerando a identidade étnico-racial, religiosa, de gênero e de sexualidade e os(as) estudantes público-alvo de uma educação inclusiva para o Ensino Fundamental.

% matched lemmas: brincadeira, brincar, criança, cuidado, cuidar, dia, doméstico, infância, instituição, lúdico, partilha, pequeno, sugestão
\end{textsample}
